\section{A nonuniform EI intervention on degree dispersion}
\label{sec:nonuniform-degree}

This test changes count fluctuations while preserving the marginal mean and
variance of every synapse and the population fractions. It directly tests the
conditional construction in a network where E and I have different local
laws. It is distinct from the earlier uniform-variance degree campaign.

\subsection{Microscopic control and exact finite-pool identities}

Let a postsynaptic neuron of type \(a\) sample a uniformly chosen subset of a
presynaptic population of size \(M=N_b=f_bN\). Write \(A_j\in\{0,1\}\) for its
edge indicators, \(K=\sum_{j=1}^M A_j\), and
\[
 E K=Mp,\qquad
 \operatorname{Var}K=\chi_{\rm eff}Mp(1-p).
\]
The subset is uniform conditional on \(K\), and distinct receiving rows are
drawn independently. Consequently
\[
 E A_j=\frac{E K}{M}=p,\qquad
 \operatorname{Var}A_j=p(1-p).
\]
For \(j\ne k\), exchangeability gives
\[
 \operatorname{Var}K
 =M p(1-p)+M(M-1)\operatorname{Cov}(A_j,A_k),
\]
hence
\[
 \operatorname{Cov}(A_j,A_k)
 =\frac{(\chi_{\rm eff}-1)p(1-p)}{M-1}.
\]
For \(J_{ij}=t_{ab}A_j/\sqrt N\), choose
\[
 p_{ab}=\frac{(J^0_{ab})^2}{(J^0_{ab})^2+G_{ab}},
 \qquad t_{ab}=\frac{J^0_{ab}}{p_{ab}}.
\]
Then \(t_{ab}\) has the presynaptic Dale sign and
\begin{align*}
 E J_{ij}&=\frac{t_{ab}p_{ab}}{\sqrt N}
          =\frac{J^0_{ab}}{\sqrt N},\\
 \operatorname{Var}J_{ij}
 &=\frac{t_{ab}^2p_{ab}(1-p_{ab})}{N}
   =\frac{G_{ab}}{N}.
\end{align*}
These marginal identities hold at finite \(M\), including nonzero rounding
dispersion. The within-row covariance matrix of the centered weights is
\[
 S_{ab}=\frac{G_{ab}}N
 \left[
 \frac{M-\chi_{\rm eff}}{M-1}I_M+
 \frac{\chi_{\rm eff}-1}{M-1}{\boldsymbol 1}{\boldsymbol 1}^{\mathsf T}
 \right].
\]
Preserving single-edge moments therefore does not preserve joint edge
independence. The low-rank term must be controlled when using a two-site
closure.

At \(\chi=1\), the degree construction recovers the full independent
Bernoulli law, not just its first two moments. For any binary row \(a\) with
\(k\) active entries,
\[
 P(A=a)=\frac{P(K=k)}{\binom Mk}
       =p^k(1-p)^{M-k}.
\]
Thus independent rows give exactly the same connectivity ensemble as the
Bernoulli baseline. Fresh seed labels with the same dynamics and estimator
can be combined in a separately labeled pooled diagnostic; the original
prospective comparisons remain separately identifiable.

For fixed source means \(m_j\), set
\(\bar m=M^{-1}\sum_jm_j\) and
\(V_m=M^{-1}\sum_j(m_j-\bar m)^2\). The exact cavity-row contraction is
\begin{align*}
 \operatorname{Var}\left(\sum_j\delta J_{ij}m_j\right)
 &=m^{\mathsf T}S_{ab}m\\
 &=\frac{G_{ab}}N\left[
 \frac{M-\chi_{\rm eff}}{M-1}M(V_m+\bar m^2)
 +\frac{\chi_{\rm eff}-1}{M-1}M^2\bar m^2\right]\\
 &=G_{ab}f_b\left[
 \frac{M-\chi_{\rm eff}}{M-1}V_m+\chi_{\rm eff}\bar m^2
 \right].
\end{align*}
The last simplification uses
\(M-\chi_{\rm eff}+M(\chi_{\rm eff}-1)=\chi_{\rm eff}(M-1)\).
Thus the usual static equation is its leading dense-network limit, rather
than an exact finite-network identity.

The count generator records the expected dispersion, not merely the
requested parameter. For \(0\le\chi\le1\), define
\(L=\operatorname{round}(\chi M)\),
\(b=(M-L)p\), and \(\rho=b-\lfloor b\rfloor\). It draws
\[
 K=\lfloor b\rfloor+B_\rho+B_{L,p},
\]
with independent Bernoulli and binomial terms. Therefore
\[
 E K=Mp,\qquad
 \chi_{\rm eff}
 =\frac{L p(1-p)+\rho(1-\rho)}{Mp(1-p)}.
\]
In particular, requested \(\chi=0\) gives exactly fixed counts only when
\(Mp\) is an integer. For \(1<\chi<M\), the beta-binomial construction uses
concentration \(c=(M-\chi)/(\chi-1)\), so that
\[
 \operatorname{Var}K
 =Mp(1-p)\frac{M+c}{1+c}
 =\chi Mp(1-p).
\]

\subsection{Conditional one-site prediction}

Under the same asynchronous dense-network cavity assumptions as the main
construction, the leading quenched input variance is
\[
 \Delta_a=\sum_bG_{ab}f_b
 \left[\operatorname{Var}_z m_b(z)+\chi_{ab}\bar m_b^2\right].
\]
The mean-rate term has a direct microscopic interpretation. At bounded
nonzero dispersion, \(K-Mp=O(\sqrt M)\). Multiplying that count deviation
by a synapse of size \(O(N^{-1/2})\) and a nonzero mean rate produces an
\(O(1)\) fixed input offset. Relative degree fluctuations can therefore
vanish while the operating-point spread persists. This does not require
the conditional temporal input variance to fail to concentrate.
Write \(\mu_a(z)=h_a+\sqrt{\Delta_a}z\),
\(q_a=K_a(0)\), and
\[
 m_a(z)=2\Phi\!\left(\frac{\mu_a(z)}{\sqrt{1+q_a}}\right),
 \qquad
 h_a=\sqrt{1+\Delta_a+q_a}\,
       \Phi^{-1}\!\left(\frac{\bar m_a}{2}\right).
\]
The last relation fixes the leading ensemble mean; no measured network mean
is inserted. Let
\[
 c_a(z,\tau)=
 E[\phi(\mu_a(z)+u_a(t+\tau))\phi(\mu_a(z)+u_a(t))]
 -m_a(z)^2,
\]
where \(u_a\) is Gaussian with covariance \(K_a\). The kernel equation is
\[
 K_a(\tau)-K_a''(\tau)
 =\sum_bG_{ab}f_b E_z c_b(z,\tau),
 \qquad K_a'(0)=0,\quad K_a(\infty)=0.
\]
The time-averaged conditional response gain is
\[
 \alpha_a(z)=
 \sqrt{\frac{2}{\pi(1+q_a)}}\,
 \exp\!\left[-\frac{\mu_a(z)^2}{2(1+q_a)}\right].
\]
The whole kernel and the conditional distribution are recomputed as
\(\chi\) changes; keeping only mean gains fixed would be a different
experiment.

For a temporally centered population covariance matrix \(C_b(\tau)\), the
finite-pool dynamic row contraction is exactly
\[
 \operatorname{tr}[S_{ab}C_b(\tau)]
 =\frac{G_{ab}}N\left[
 \frac{M-\chi_{\rm eff}}{M-1}\operatorname{tr}C_b(\tau)
 +\frac{\chi_{\rm eff}-1}{M-1}
 {\boldsymbol 1}^{\mathsf T}C_b(\tau){\boldsymbol 1}\right].
\]
The second term is subleading if the population-mean covariance vanishes
sufficiently quickly. Stable coherent feedback and a self-averaging
asynchronous state are substantive conditions, not consequences of matching
the single-edge moments. The paired-row contraction and its directly
measured low-rank diagnostics are treated separately in the finite-pool audit.

\subsection{Geometry and observables}

Let \(v_a(z)=c_a(z,0)\),
\(A_a=E_z\alpha_a(z)^2\), and
\(Q_a(\omega_1,\omega_2)=
E_z[c_a(z,\omega_1)c_a(z,\omega_2)]\).
With \(F=\operatorname{diag}(f)\),
\(H(\omega)=(1+i\omega)^{-1}\), and
\[
 R(\omega_1,\omega_2)=
 \left[I-\operatorname{diag}(A)H(\omega_1)H(\omega_2)GF\right]^{-1},
\]
the conditional pair closure gives the integrated matrix
\[
 X=\iint R\operatorname{diag}(Q/f)R^\dagger
       \frac{d\omega_1d\omega_2}{(2\pi)^2}.
\]
Our implementation stores spectral masses already containing the integration
measures. Its frequency sums therefore receive no second quadrature factor.
Define
\[
 P_{ab}=X_{ab}-\delta_{ab}\frac{E_zv_a^2}{f_a},
 \qquad
 \mathcal B_{ab}=f_af_b P_{ab}+
                  \delta_{ab} f_aE_zv_a^2.
\]
The matrix \(\mathcal B\) is directly comparable to
\(N^{-1}\sum_{i\in a,j\in b}C_{ij}^2\).
Its two equal cross-population entries each contribute once:
\[
 \mathcal E=\mathcal B_{EE}+2\mathcal B_{EI}+\mathcal B_{II},
 \quad
 \nu=\sum_af_aE_zv_a,
 \quad
 d=\frac{\nu^2}{\mathcal E}.
\]
These are covariance energies, not participation ratios of a rectangular
E--I block. For a single population submatrix,
\[
 \frac{D_a}{N_a}
 =\frac{(E_zv_a)^2}{E_zv_a^2+f_aP_{aa}}.
\]

A useful exact decomposition within this prediction is
\[
 d=d_{\rm uncorr}
       \frac{\mathcal E_{\rm diag}}{\mathcal E},
 \qquad
 d_{\rm uncorr}=\frac{\nu^2}{\mathcal E_{\rm diag}},
 \qquad
 \mathcal E_{\rm diag}=\sum_af_aE_zv_a^2.
\]
The first factor describes the spread of individual temporal variances;
the second describes the further dimensional reduction by correlations.
Neither factor is a measure of computational task performance.

\subsection{Specified nonuniform control and deterministic checks}

The control uses
\[
 f=(0.8,0.2),\qquad
 J^0=\begin{pmatrix}1&-6\\2&-10\end{pmatrix},\qquad
 G=\begin{pmatrix}16&9\\25&12\end{pmatrix},\qquad
 \bar m_E=\bar m_I=0.3.
\]
These quantities remain fixed while all four degree-dispersion entries take
the common value \(\chi\). The external input is
\[
 I_a^{\rm total}=
 -\sqrt N\sum_bJ^0_{ab}f_b\bar m_b+h_a(\chi).
\]
The recomputed \(h_a(\chi)\) maintains the intended leading operating point.
It is fixed before each simulation and is never fit to the observed PR.
The \(O(\sqrt N)\) balance drive is fixed along this path, whereas the
\(O(1)\) offset changes with \(\chi\). This is not an intervention with every
external input held fixed. In the balanced limit the large drive already
fixes the mean rates; setting its finite offset to the limiting \(h_a(\chi)\)
also aims to reduce finite-gain mean offsets without fitting them.

\begin{center}
\begin{tabular}{r|rr|rr|rrr}
\(\chi\)&\(\Delta_E\)&\(\Delta_I\)&\(q_E\)&\(q_I\)&
\(d_{\rm uncorr}\)&\(\mathcal E_{\rm diag}/\mathcal E\)&\(d\)\\
\hline
0&0&0&3.24483&4.96691&0.99637&0.17452&0.17389\\
0.5&2.14132&3.28237&2.48054&3.79709&0.57194&0.09781&0.05594\\
1&3.70336&5.67745&1.93951&2.96866&0.45827&0.07490&0.03433\\
2&6.13630&9.40862&1.24928&1.91174&0.34047&0.04937&0.01681
\end{tabular}
\end{center}

Refining the one-site time grid from \(0.08\) to \(0.04\), the horizon from
\(100\) to \(150\), the Gaussian quadrature from \(96\) to \(160\) points, and
the spectral cutoff from \(12\) to \(20\) changes \(d\) by at most
\(3.3\times10^{-8}\) in relative value at these four points. This checks
numerical implementation of the conditional closure; it does not validate
the closure against networks.

At \(N=1000\), the exact expected rounding dispersions for requested
\(\chi=0\) are
\[
 \begin{pmatrix}0.00125&0\\0.002375&0.0128\end{pmatrix}.
\]
Re-solving the conditional equations with these values while reimposing the
target means changes \(d\) by about \(-0.74\%\) and either component of \(h\)
by less than \(0.003\). This is a rounding sensitivity diagnostic, not an
exact prediction for the already specified drive, because it redefines
\(h\). It is small compared with the predicted 10.34-fold endpoint change.

All network comparisons use fresh prespecified connectivity draws and retain
every outcome. The mean-free covariance-energy estimator is compared at fixed
block duration and checked against longer records. Its independence and
stationarity requirements remain explicit. A finite-record ratio is not
made unbiased merely by correcting its denominator.

\subsection{Regular static susceptibility on this control path}

The fixed-count endpoint also admits a direct local check of the degree
mechanism. At \(\chi=0\), the leading static law has \(\Delta=0\), while the
temporal variance remains positive. Expanding the conditional mean in its
small quenched offset gives
\[
 m_b(z)=\bar m_b+\alpha_{b,0}\sqrt{\Delta_b}z+O(\Delta_b),
 \qquad
 \operatorname{Var}_z m_b(z)=A_{b,0}\Delta_b+O(\Delta_b^2).
\]
The mean is held fixed, so changes of \(h_b\) and \(q_b\) contribute only at
higher order to this variance coefficient. For a common small dispersion,
the static equation therefore reads
\[
 \Delta=W\operatorname{diag}(A_0)\Delta+
          \chi W\bar m^2+O(\|\Delta\|^2+\chi\|\Delta\|),
 \qquad W=GF.
\]
If \(\rho(W\operatorname{diag}A_0)<1\), its derivative is
\[
 \left.\frac{d\Delta}{d\chi}\right|_{0}
 =[I-W\operatorname{diag}(A_0)]^{-1}W\bar m^2.
\]
For the specified profile, the direct degree source and full susceptibility are
\[
 W\bar m^2=(1.314,2.016)^{\mathsf T},\qquad
 \Delta'(0)=(5.390674,8.261608)^{\mathsf T}.
\]
The positive response eigenvalue is
\(0.756211\). Static recurrent feedback amplifies the direct count source by
about four at this endpoint. Numerical solutions at
\(\chi=10^{-3},10^{-4},10^{-5}\) approach the derivative with maximum
relative errors \(6.47\times10^{-4},6.48\times10^{-5},6.48\times10^{-6}\).
This is a regular, away-from-onset susceptibility. It is not a claim that the
same linear expansion stays uniform at a critical point.
